% Part: counterfactuals
% Chapter: introduction
% Section: counterfactuals

\documentclass[../../../include/open-logic-section]{subfiles}

\begin{document}

\olfileid{cnt}{int}{cnt}

\olsection{في الشرطيات المخالفة للواقع}

ومن أشهر صيغ «إذا \dots، فإن \dots» الإنجليزية وأهمها ما يُبنى بصيغة
الماضي الالتزامية من فعل \emph{يكون}: «لو كان الأمر أن \dots، لكان
الأمر أن \dots». ولما كان مقدّم هذا الضرب من الشرط كاذبًا في العادة،
أي مخالفًا للواقع، سُمّيت تراكيبه \emph{شرطيات مخالفة للواقع}؛ وتسمى
أيضًا \emph{شرطيات التزامية} لاستعمالها تلك الصيغة من فعل \emph{يكون}.
وهي غير الشرطيات \emph{الخبرية} التي صورتها «إذا كان الأمر أن \dots،
فإن الأمر أن \dots»؛ بل تختلف عنها في شروط الصدق. ومثال آدامز المشهور
يكشف الفرق:
\begin{quote}
  إذا لم يقتل أوزوالد كينيدي، فقد قتله شخص آخر.
  
  لو لم يكن أوزوالد قد قتل كينيدي، لكان شخص آخر قد قتله.
\end{quote}
فالأولى خبرية، والثانية مخالفة للواقع. وصدق الأولى بيّن، لأن المعلوم أن
الرئيس جون ف. كينيدي قتله \emph{شخص ما}؛ فإن لم يكن القاتل، خلافًا
لتقرير وارن، لي هارفي أوزوالد، كان غيره. وأما الثانية فدعواها أبعد:
مؤداها أنه لو لم يقع قتل أوزوالد لكينيدي، بأن مُنع إطلاق النار في دالاس
أو أخفق، لسار التاريخ بعد ذلك إلى نجاح محاولة اغتيال أخرى. ولا تصدق
إلا إذا كانت قوى نافذة قد دبرت لضمان موت جون ف. كينيدي، كما يذهب كثير
من أصحاب نظريات المؤامرة في اغتياله.

ولا يزال الخلاف قائمًا في صحة صوغ الشرط \emph{الخبري} بالشرط المادي،
ولا سيما في إمكان «تفسير» مفارقاته تفسيرًا يوافق كونه محددًا لشروط صدق
الشرطيات الخبرية الإنجليزية. أما قصور الشرط المادي عن التمثيل الصحيح
للشرطيات المخالفة للواقع فمتفق عليه. والسبب أن كذب مقدماتها في الغالب
لا يجعل كل شرط مخالف للواقع ذي مقدّم كاذب صادقًا؛ وشرط آدامز شاهد، إذا
قبلنا تقرير وارن ونفينا وجود مؤامرة لاغتيال كينيدي.

وللشرطيات المخالفة للواقع منزلة ظاهرة في الاستدلال السببي، إذ يُعبّر
بها عن العلل وآثارها. فحك عود الثقاب سبب اشتعاله، ويقال: «لو حُك هذا
العود، لاشتعل». ولا يؤدي الشرط المادي، ولا الشرط الخبري في الجملة، هذا
المعنى؛ لأن «يُحك عود الثقاب~$\lif$ يشتعل عود الثقاب» يصدق إذا لم
يُحك العود أصلًا، كيفما كان سيؤول أمره لو حُك. بل يصدق حينئذ أيضًا
«يُحك عود الثقاب~$\lif$ يتحول عود الثقاب إلى باقة زهور»، مع القطع بأنه
لو حُك لم ينقلب زهورًا.

ولا اتفاق إلى الآن على المنطق الصحيح لهذه الشرطيات. ومن أشهر ما قُدّم
تحليل ستالنيكر ولويس: فالشرط «لو كان الأمر أن~$\OLId{latin-looped}{S}$، لكان الأمر أن~$\OLId{latin-looped}{T}$»
يصدق، في مذهبهما، إذا وفقط إذا صدقت~$\OLId{latin-looped}{T}$ في أقرب حالة مخالفة للواقع، أو
«عالم ممكن»، إلى العالم الفعلي من الحالات التي تصدق فيها~$\OLId{latin-looped}{S}$. ويسمى
التحليل «وجوديًا» لرده الأمر إلى أنطولوجيا العوالم الممكنة. وترد تحليلات
أخرى إلى الاحتمال الشرطي أو إلى نظريات مراجعة الاعتقاد، فتعددت المنطقيات
المقترحة. وليس عند لويس وستالنيكر نفسيهما منطق واحد؛ فتحليلاهما وإن
تقاربا في الرجوع إلى أقرب العوالم الممكنة، تختلف فروضهما فتختلف معها
الاستدلالات الصحيحة.

\end{document}



